%%%%%
%
\section{Complete\_X\_avec\_Y}

\itemstyle{Op\'erateur~:} {\sf Complete\_X\_avec\_Y}
\index{Complete\_X\_avec\_Y}
\label{Complete_X_avec_Y}

%%%%%OBJET%%%%
\itemstyle{Objet~:}
Complete un vecteur X avec un vecteur Y. Cet op�rateur cr�e donc
un troisi�me vecteur Z qui contient des valeurs du vecteur X 
et des valeurs du vecteur Y lorsque celles de X sont manquantes.

Un flag en sortie indique de quel vecteur sont issues les
valeurs du vecteur Z en sortie.

%%%%%ARGUMENTS%%%%
\bigskip
%%%%%ENTREES%%%%
\begin{tabarg}
\multicolumn{3}{l}{\titretab {Entr\'ees}}\\
\hline
$X$ & Vecteur, \texttt{float} & Vecteur sur lequel on appliquera le test (valeur manquante?)\\
$Y$ & Vecteur, \texttt{float} & Vecteur de remplacement, lorsque qu'un 
�chantillon de $X$ est manquant(m�me unit� que $X$).\\
$valmnqX$ & Coeff., \texttt{float} & Codage de la valeur manquante du vecteur $X$
(n'est utile que pour les vecteurs bugu�s dont l'attribut misval est 
absent, ou ne correspond pas � la r�alit�).\\
\end{tabarg}

\medskip
%%%%%SORTIES%%%%
\begin{tabarg}
\multicolumn{3}{l}{\titretab {Sorties}}\\
\hline
$flag$ & Vecteur, \texttt{float} & Flag pour indiquer la provenance de la valeur de $Z$\\
$Z$ & Vecteur, \texttt{float} & Vecteur coomposite (m�me unit� que $X$)\
\end{tabarg}


%%%%%FORMULE%%%%
\itemstyle{Formule}

%%Complete_X_avec_Y
1- Recherche des �ventuelles valeurs manquantes dans le vecteur $X$ \\
2- Remplacement si besoin par des valeurs du vecteur $Y$

En sortie : le flag permet de conna�tre l'origine des valeurs :\\
\hspace*{2cm} 1 : si la valeur du vecteur $Z$ provient du vecteur d'origine $X$\\
\hspace*{2cm} 2 : si la valeur du vecteur $Z$ provient du vecteur de remplacement $Y$\\



%%%%%COMMENTAIRE%%%%
\itemstyle{Commentaire}

Ne pas oublier que les vecteurs $X$ et $Y$ doivent avoir le m�me 
nombre d'�chantillons.

L'utilisation de cet op�rateur est proche de l'utilisation de
l'op�rateur select\_sur\_vect (\ref{select_sur_vect}) :  ce dernier 
compl�te un param�tre 
lorsque ses valeurs sont en dehors de bornes fix�es.



%%%%%AUTEUR%%%%
\itemstyle{Auteur~:}
CNRM/GMEI/TRAMM

%%%%%REFERENCE%%%%
\itemstyle{R\'ef\'erences~:}

